%=====================================================================
\chapter{Intelligent Agents and Task Environments}\label{ch:agents}
%=====================================================================
\locator{1}{Part I --- Foundations of Intelligent Systems}

\begin{outcomes}
By the end of this chapter you will be able to:
\begin{tight}
  \item \textbf{Define} an agent in terms of percepts, actions, an agent function
        and an agent program, and \textbf{explain} why the distinction between the
        last two matters.
  \item \textbf{Specify} any task for an agent as a \textbf{PEAS} description ---
        performance measure, environment, actuators, sensors --- before writing a
        line of code.
  \item \textbf{Classify} a task environment on the six standard properties, and
        \textbf{predict} which agent architectures the classification rules out.
  \item \textbf{Distinguish} the five kinds of agent program and \textbf{choose}
        the simplest one that can do a given job.
  \item \textbf{Explain} rationality as doing the best that can be expected given
        the percepts so far, and \textbf{separate} it from omniscience and from
        perfection.
\end{tight}
\end{outcomes}

\begin{prereq}
\chref{introduction} (the definition of an agent, the acting-rationally quadrant,
the five-point checklist). The agent vocabulary set up here is used by every later
chapter: search is what a goal-based agent does to choose its action
(\chref{uninformed}), and utility appears again as expected utility in
\chref{uncertainty}.
\end{prereq}

\begin{keyterms}
agent $\bullet$ percept $\bullet$ percept sequence $\bullet$ actuator $\bullet$
sensor $\bullet$ agent function $\bullet$ agent program $\bullet$ PEAS $\bullet$
performance measure $\bullet$ rational agent $\bullet$ fully and partially
observable $\bullet$ deterministic and stochastic $\bullet$ episodic and
sequential $\bullet$ static and dynamic $\bullet$ discrete and continuous
$\bullet$ single- and multi-agent $\bullet$ simple reflex agent $\bullet$
model-based reflex agent $\bullet$ goal-based agent $\bullet$ utility-based agent
$\bullet$ state estimation
\end{keyterms}

\section{Agent, Environment, Percept, Action}

Everything in this course is an agent doing one thing: taking in what it can see
and choosing what to do. \dsref{agent-loop} is the picture, and it is worth
looking at carefully, because four of its labels are technical terms that get used
precisely from here on.

\dsfig{%
\begin{tikzpicture}[font=\small,
  ag/.style={rounded corners=4pt, draw=primary, line width=1.4pt, fill=primary!8,
             minimum width=5.4cm, minimum height=3.5cm},
  env/.style={rounded corners=4pt, draw=goldhl, line width=1.4pt, fill=goldhl!8,
              minimum width=3.6cm, minimum height=3.5cm, align=center,
              text=goldhl!50!black},
  slot/.style={rounded corners=3pt, fill=white, draw=primary!45, line width=0.8pt,
               font=\scriptsize\bfseries, text=primary!70!black, align=center,
               minimum width=4.4cm, minimum height=0.72cm},
  a/.style={-{Stealth[length=2.6mm]}, line width=1.2pt}]

\node[ag] (agent) at (0,0) {};
\node[font=\scriptsize\bfseries, text=primary] at (0,1.52) {AGENT};
\node[slot] (prog) at (0,0.55) {agent program};
\node[slot] (fn)   at (0,-0.45)
  {\emph{agent function}\\[-1pt]{\tiny percept sequence $\mapsto$ action}};
\node[font=\tiny\itshape, text=neutral, text width=4.6cm, align=center]
  at (0,-1.3) {the program \emph{implements} the function\\ in bounded time and
  memory};

% Every \\ in a node stays at the top level of the node text: a \\ inside a
% nested {...} group breaks the alignment (see the README's risks section).
\node[env, right=3.4cm of agent] (world) {ENVIRONMENT\\[3pt]
  {\scriptsize the three hundred}\\
  {\scriptsize open tickets, the}\\
  {\scriptsize engineers on shift,}\\
  {\scriptsize the clock}};

\draw[a, secondary] (world.north west) -- ++(-1.0,0)
  node[midway, above, font=\scriptsize, text=secondary]
  {\textbf{percept}} -- (agent.north east);
\draw[a, accent] (agent.south east) -- ++(1.0,0)
  node[midway, below, font=\scriptsize, text=accent]
  {\textbf{action}} -- (world.south west);

\node[font=\tiny, text=secondary, align=center] at (3.62,0.46)
  {\emph{via} sensors};
\node[font=\tiny, text=accent, align=center] at (3.62,-0.5)
  {\emph{via} actuators};

\node[rounded corners=3pt, fill=greenhl!10, draw=greenhl, line width=0.9pt,
      text=greenhl!45!black, font=\scriptsize, text width=12.2cm,
      align=flush left, inner sep=6pt] at (1.7,-3.1)
  {\textbf{The escalation adviser, in these terms.} \emph{Percepts:} the ticket
   queue, each ticket's age, severity and service level agreement, the shift
   roster, each engineer's current load. \emph{Actions:} escalate a ticket to a
   named engineer, or leave it. \emph{Agent function:} the mapping from any
   possible queue-and-roster situation to a set of escalations. \emph{Agent
   program:} the two hundred lines of Chapter~13 that compute one of those
   mappings tonight, in under a second.};
\end{tikzpicture}}%
{The agent and its environment. The \emph{agent function} is the mathematical
object --- what the agent ought to do in every situation it could ever be in. The
\emph{agent program} is the code that runs on the machine.}{agent-loop}

\begin{definitionbox}[Agent function and agent program]
The \term{agent function} maps every possible \term{percept sequence} to an
action:
\[
  f:\; \mathcal{P}^{*} \longrightarrow \mathcal{A}
\]
where $\mathcal{P}^{*}$ is the set of all finite sequences of percepts and
$\mathcal{A}$ the set of available actions. It is an abstract object, usually
infinite, and it is what we are trying to specify.

The \term{agent program} is a concrete implementation of $f$ that runs inside the
agent, with finite memory and a deadline.
\end{definitionbox}

The distinction looks pedantic and is not. The agent function for the escalation
problem could be written as an enormous table: every possible state of the queue in
the left column, the right set of escalations in the right column. That table would
be a complete, correct specification --- and with three hundred tickets it has more
rows than there are atoms in the observable universe. So the interesting work of AI
is never ``what is the right answer in every case''; it is ``what program computes
a good enough answer in the time available''. Every method in this book is an
answer to that second question.

\begin{conceptbox}
Notice also that the agent function takes the whole percept \emph{sequence}, not
the latest percept. That single choice is the difference between the first two
agent architectures in Section~\ref{sec:agent-kinds}. An agent that reacts only to
what it can see right now cannot escalate ``the ticket that has been quiet
suspiciously long'', because quietness is not visible in a snapshot --- it is a
property of the history.
\end{conceptbox}

\section{PEAS: Specify the Task Before You Build the Agent}

Before any code, write down four things. The mnemonic is \term{PEAS}.

\begin{itemize}[leftmargin=2.4em, itemsep=3pt, topsep=4pt]
  \item \textbf{P}erformance measure --- what counts as success, stated so that two
        people would score the same run identically.
  \item \textbf{E}nvironment --- what the agent operates in, including the parts it
        cannot see.
  \item \textbf{A}ctuators --- what the agent can do.
  \item \textbf{S}ensors --- what the agent can perceive.
\end{itemize}

The performance measure comes first and is the one people get wrong. It must be a
measure of the \emph{environment's} state, not of the agent's behaviour. An
escalation adviser scored on ``number of escalations made'' will escalate
everything. Scored on ``breached service level agreements, plus a penalty for each
escalation that turned out to be unnecessary'', it has to make a judgement --- and
now the problem is interesting.

\smallskip
\noindent\begin{longtable}{@{}L{2.9cm}L{3.5cm}L{3.5cm}L{4.1cm}@{}}
\toprule
\textbf{Task} & \textbf{Performance measure} & \textbf{Environment} &
\textbf{Actuators / Sensors} \\
\midrule
\endfirsthead
\toprule
\textbf{Task} & \textbf{Performance measure} & \textbf{Environment} &
\textbf{Actuators / Sensors} \\
\midrule
\endhead
Escalation adviser \newline \emph{(Ch.\ 13)} &
Breached service level agreements; penalty per unnecessary escalation; on-call
hours consumed &
Ticket queue, shift roster, the clock; other engineers acting independently &
\emph{Act:} escalate to a named engineer, defer, close as duplicate.
\emph{Sense:} ticket fields, roster, current loads \\
Sprint scheduler \newline \emph{(Ch.\ 7)} &
Stories placed; zero constraint violations; developer load balance &
The backlog, three sprint slots, capacity and precedence rules &
\emph{Act:} assign a story to a slot, unassign.
\emph{Sense:} the full backlog and all constraints \\
Release planner \newline \emph{(Ch.\ 9)} &
Plan length; whether the goal state is actually reached &
Code, review, test and deploy states of each component &
\emph{Act:} apply an operator (code, review, test, deploy, roll back).
\emph{Sense:} the current state of every component \\
Contract negotiator \newline \emph{(Ch.\ 8)} &
Value of the agreed slot to us, given that they are maximising theirs &
Two parties, alternating offers, a deadline &
\emph{Act:} offer, accept, walk away.
\emph{Sense:} their last offer and the history \\
Defect diagnoser \newline \emph{(Ch.\ 15)} &
Probability mass placed on the true cause; cost of wrong remediation &
A live system; causes not directly observable &
\emph{Act:} report a posterior over causes, order one test.
\emph{Sense:} symptoms, logs, test outcomes \\
\bottomrule
\caption{PEAS descriptions for five tasks from one project office. Each row names
the chapter that builds the agent.}\label{tab:peas}
\end{longtable}

\section{Six Properties of a Task Environment}

Once the task is specified, classify the environment. The classification is not
academic bookkeeping: each property rules some architectures in and others out, and
getting it wrong is how projects fail in month four.

\smallskip
\noindent\begin{longtable}{@{}L{3.3cm}L{5.5cm}L{6.2cm}@{}}
\toprule
\textbf{Property} & \textbf{The question} & \textbf{Where it bites} \\
\midrule
\endfirsthead
\toprule
\textbf{Property} & \textbf{The question} & \textbf{Where it bites} \\
\midrule
\endhead
Fully \emph{vs.} partially observable &
Do the sensors give the complete state relevant to the choice? &
Partially observable means the agent must \emph{maintain} state --- an internal
model updated over time. This is the single most consequential property: it is the
whole difference between the two reflex agents of
Section~\ref{sec:agent-kinds}. \\
Deterministic \emph{vs.} stochastic &
Does the current state plus the action fix the next state? &
Deterministic lets you search for a plan once and execute it
(\chref{planning}). Stochastic forces probability (\chref{uncertainty}) or
re-planning after every step. \\
Episodic \emph{vs.} sequential &
Is each decision independent, or does it constrain later ones? &
Episodic problems can be handled one case at a time. Sequential ones need
lookahead: escalating your only free engineer at 21:00 is what makes the 23:00
decision hard. \\
Static \emph{vs.} dynamic &
Can the world change while the agent is deliberating? &
Dynamic environments put a deadline on thinking, which is why
\chref{adversarial} spends as much effort on cutting search off as on searching. \\
Discrete \emph{vs.} continuous &
Are states, time and actions countable or real-valued? &
Discrete gives you the graph of \chref{uninformed}. Continuous needs the local
search of \chref{local}. \\
Single- \emph{vs.} multi-agent &
Is anything else in the environment choosing its actions? &
If another party's best move depends on yours, you need
\chref{adversarial}; treating a negotiator as weather is the classic error. \\
\bottomrule
\caption{The six properties, and the design decision each one forces.}
\label{tab:envprops}
\end{longtable}

\begin{alertbox}[Classify honestly, not conveniently]
The escalation problem is \emph{partially observable} (you cannot see which
tickets are about to become urgent), \emph{stochastic} (whether an escalation
helps depends on the engineer), \emph{sequential}, \emph{dynamic} (tickets arrive
while you think), \emph{discrete} and \emph{multi-agent} (engineers make their own
choices). That is the hard corner of every one of the six.

The standard mistake is to notice this, find it discouraging, and quietly build
for the easy corner anyway --- a fully observable, deterministic, episodic model
--- because that is what the tutorial did. The resulting system works in
demonstrations and fails on a Friday night. The honest response is to
\emph{simplify deliberately and write the simplification down}: this is exactly
what \chref{csp} does when it treats a sprint as static for the duration of the
solve, and says so.
\end{alertbox}

\section{Five Kinds of Agent Program}
\label{sec:agent-kinds}

\dsref{agent-kinds} shows the five architectures in increasing order of what they
keep inside. The professional skill is to choose the \emph{simplest} one that can
do the job, because each step up costs code, state and ways to be wrong.

\dsfig{%
\begin{tikzpicture}[font=\scriptsize,
  box/.style={rounded corners=3pt, draw=#1, line width=1.1pt, fill=#1!8,
              text width=2.55cm, align=flush left, inner sep=5pt,
              text=#1!50!black, minimum height=3.2cm},
  hd/.style={rounded corners=2pt, fill=#1, text=white, font=\scriptsize\bfseries,
             inner sep=2.5pt, text width=2.4cm, align=center},
  a/.style={-{Stealth[length=2mm]}, neutral!70, line width=0.9pt}]

\node[box=primary] (k1) at (0,0)
  {\\[2pt]Condition--action rules on the \emph{current} percept only.\\[3pt]
   \emph{Keeps:} nothing.\\[3pt]
   {\itshape Escalate any ticket already past its deadline.}\\[3pt]
   Fails the moment the right action depends on history.};

\node[box=secondary, right=0.42cm of k1] (k2)
  {\\[2pt]Same rules, plus an internal \emph{model} of the unobserved state,
   updated each cycle.\\[3pt]
   \emph{Keeps:} an estimate of the world.\\[3pt]
   {\itshape Track each engineer's load, which no percept reports directly.}};

\node[box=greenhl, right=0.42cm of k2] (k3)
  {\\[2pt]Holds an explicit \emph{goal} and searches for an action sequence that
   reaches it.\\[3pt]
   \emph{Keeps:} model + goal.\\[3pt]
   {\itshape Reach a state where no service level agreement is breached.}\\[3pt]
   This is Chapters~4--9.};

\node[box=tealhl, right=0.42cm of k3] (k4)
  {\\[2pt]Replaces the yes/no goal with a \emph{utility} over outcomes, and
   maximises expected utility.\\[3pt]
   \emph{Keeps:} model + utility.\\[3pt]
   {\itshape Trade two minor breaches against one major one.}\\[3pt]
   This is Chapter~14.};

\node[box=purplehl, right=0.42cm of k4] (k5)
  {\\[2pt]Any of the above, plus a component that \emph{improves} it from
   experience.\\[3pt]
   \emph{Keeps:} everything + a critic.\\[3pt]
   {\itshape Learn which escalations were worth making.}\\[3pt]
   Chapter~17, then \emph{Machine Learning}.};

\node[hd=primary]   at (k1.north) {1. SIMPLE REFLEX};
\node[hd=secondary] at (k2.north) {2. MODEL-BASED REFLEX};
\node[hd=greenhl]   at (k3.north) {3. GOAL-BASED};
\node[hd=tealhl]    at (k4.north) {4. UTILITY-BASED};
\node[hd=purplehl]  at (k5.north) {5. LEARNING};

\foreach \x/\y in {k1/k2,k2/k3,k3/k4,k4/k5}{\draw[a] (\x.east) -- (\y.west);}

\draw[decorate, decoration={brace, amplitude=5pt, mirror}, neutral, line width=0.9pt]
  ($(k1.south west)+(0,-0.15)$) -- ($(k4.south east)+(0,-0.15)$)
  node[midway, below=6pt, font=\scriptsize\itshape, text=neutral]
  {the four architectures this course builds};
\end{tikzpicture}}%
{Five agent architectures. Each keeps more state than the one before it. Choose the
leftmost one that can do the job.}{agent-kinds}

\section{Rationality Is Not Omniscience}

A \term{rational agent}, for every percept sequence, selects the action expected to
maximise its performance measure, given the evidence the percept sequence provides
and whatever built-in knowledge it has.

Three things this definition deliberately excludes.

It does not require \emph{omniscience}. An agent that escalates ticket~114 on
excellent evidence, and turns out to have been wrong because the customer had
already withdrawn the complaint by e-mail that nobody had read, was rational and
unlucky. Rationality is a property of the decision procedure, judged against the
information available, not a property of the outcome. This distinction is what
lets you review a bad night fairly.

It does not require \emph{perfection}. Perfection maximises actual performance;
rationality maximises expected performance. Only the second is achievable, and
demanding the first produces agents that stall.

And it does not forbid \emph{information gathering}. If the agent can act to
improve its own percepts --- run a test, ask the customer, look at the log --- then
doing so is often the rational action. \chref{uncertainty} puts a number on
exactly how much such an act is worth.

\begin{worked}[PEAS and a full classification for one night's triage]
\textbf{The task.} At 21:00 the queue holds three hundred open tickets. One
engineer, Ayesha, is on call until 07:00. Tickets breach at different times.
Escalating wakes Ayesha, which costs something, and there is a limit to how often
it can be done in a night.

\smallskip
\textbf{PEAS.}
\begin{tight}
  \item \emph{Performance measure:} $-10$ per breached major service level
        agreement, $-3$ per breached minor one, $-1$ per escalation that turned
        out unnecessary, $-4$ per hour of Ayesha's sleep lost. Note it is a
        measure of the world's state at 07:00, not of how busy the agent was.
  \item \emph{Environment:} the queue, the roster, the clock, the customers
        (who are also acting), Ayesha (who is also acting).
  \item \emph{Actuators:} escalate ticket $t$ to Ayesha; defer $t$ to the morning;
        merge $t$ into another ticket.
  \item \emph{Sensors:} every ticket's fields; Ayesha's current task; the clock.
        \emph{Not} sensed: which quiet ticket is about to escalate itself into a
        major incident.
\end{tight}

\smallskip
\textbf{Classification.} Partially observable (that last point), stochastic,
sequential (waking Ayesha at 21:00 changes what is possible at 03:00), dynamic,
discrete, multi-agent.

\smallskip
\textbf{What the classification rules out.} Partially observable kills agent~1 ---
a simple reflex agent cannot represent ``Ayesha has already been woken twice''.
Sequential kills the purely episodic treatment: you cannot score each ticket
independently and take the top five. Stochastic and multi-agent together mean the
plan will be wrong by 23:00, so any plan must be cheap enough to recompute.

\smallskip
\textbf{The design that follows.} A model-based agent that maintains Ayesha's
accumulated cost, using a goal-based search over the next few hours
(\chref{informed}), with the utility trade-off of \chref{uncertainty} to weigh a
major breach against lost sleep, recomputed every fifteen minutes. \emph{Every
one of those four choices was forced by a line of the classification}, which is
why the classification is done before the code.
\end{worked}

\begin{pitfall}
\begin{tight}
  \item \textbf{A performance measure that scores the agent instead of the world.}
        ``Number of escalations'' rewards noise; ``tickets processed'' rewards
        haste. Score the state of the environment at the end.
  \item \textbf{Confusing the agent function with the agent program.} The function
        can be an infinite table; the program has a deadline. Interview questions
        about ``the optimal policy'' are usually about the function, and
        implementation questions are always about the program.
  \item \textbf{Building a simple reflex agent for a partially observable world.}
        It will work on every test case you can construct by hand, because you can
        see the state that it cannot.
  \item \textbf{Treating another agent as weather.} If the other party's choice
        depends on yours, a stochastic model of them is not merely inaccurate --- it
        is exploitable. \chref{adversarial} is the fix.
  \item \textbf{Calling an agent irrational because the outcome was bad.} Judge the
        decision against the percepts available at the time. This is also how to
        run a blameless incident review.
\end{tight}
\end{pitfall}

%---------------------------------------------------------------------
\begin{lab}[A Reflex Agent on a Ticket Queue]
Forty tickets arrive over one night. Two agents see exactly the same percepts: a
simple reflex agent, and a model-based agent that remembers what it has already
done. The performance measure is the one from the worked example. The point of the
lab is that the second agent wins \emph{on identical information}, purely because
it keeps state.
\begin{lstlisting}[language=Python]
"""Chapter 2 lab -- a reflex agent and a model-based agent.

Two agent programs, the same percepts, the same performance measure.
The simple reflex agent cannot represent how often it has already woken
the on-call engineer, because no percept reports it. That single
missing piece of state is the whole difference.
"""

import random
from collections import namedtuple

random.seed(0)

# PEAS, as code. A percept is one ticket as the agent sees it.
Ticket = namedtuple("Ticket", "tid hour severity hours_to_breach")

WAKE_LIMIT = 3          # more than this and Ayesha is useless tomorrow
COST = {"major": 10, "minor": 3}
UNNECESSARY = 1
SLEEP = 4


def night(n=40):
    """One night's percept sequence: n tickets, 21:00 to 07:00."""
    out = []
    for tid in range(1, n + 1):
        out.append(Ticket(
            tid=tid,
            hour=21 + (tid * 10) // n,
            severity=random.choice(["major", "minor", "minor"]),
            hours_to_breach=random.choice([1, 2, 4, 8, 12]),
        ))
    return out


def simple_reflex(t, state):
    """Condition-action rules on the current percept only."""
    return t.hours_to_breach <= 4 or t.severity == "major"


def model_based(t, state):
    """Same rules, plus a model of what has already been spent.

    state["wakes"] is not in any percept: the agent maintains it.
    """
    if state["wakes"] >= WAKE_LIMIT:
        # budget gone: escalate only what will certainly breach badly
        return t.severity == "major" and t.hours_to_breach <= 2
    if t.hours_to_breach > 10:
        return False                    # cannot breach tonight; defer
    return t.hours_to_breach <= 4 or t.severity == "major"


def run(agent, tickets):
    """One night. The environment, not the agent, decides what an
    escalation achieves: past WAKE_LIMIT the engineer is too tired to
    fix anything, but the sleep is lost all the same.
    """
    state = {"wakes": 0}
    breach = sleep = waste = woke = 0
    for t in tickets:
        act = agent(t, state)
        fixed = False
        if act:
            state["wakes"] += 1
            woke += 1
            sleep += SLEEP
            if state["wakes"] <= WAKE_LIMIT:
                fixed = True            # still sharp enough to help
            if t.hours_to_breach > 10:
                waste += UNNECESSARY    # would not have breached tonight
        if not fixed and t.hours_to_breach <= 10:
            breach += COST[t.severity]
    return breach + sleep + waste, breach, sleep + waste, woke


if __name__ == "__main__":
    tickets = night()
    print("tickets:", len(tickets),
          " major:", sum(1 for t in tickets if t.severity == "major"),
          " wake limit:", WAKE_LIMIT)
    print()
    print("agent             total   breaches   on-call   escalations")
    print("-" * 58)
    rows = {}
    for name, agent in (("simple reflex", simple_reflex),
                        ("model-based  ", model_based)):
        rows[name.strip()] = run(agent, tickets)
        tot, br, oc, woke = rows[name.strip()]
        print("%-16s %6d %10d %9d %13d" % (name, tot, br, oc, woke))

    t1, b1, o1, w1 = rows["simple reflex"]
    t2, b2, o2, w2 = rows["model-based"]
    assert t2 < t1, "the model-based agent should cost less"
    assert w2 < w1, "and should wake the engineer less often"
    print()
    print("total cost  %d -> %d  (%.0f%% lower)" %
          (t1, t2, 100 * (t1 - t2) / t1))
    print("escalations %d -> %d  (%.1fx fewer)" % (w1, w2, w1 / w2))
    print()
    print("The model-based agent saw no extra percepts. It kept one")
    print("number -- how often it had already woken her -- that no")
    print("percept reports and the simple agent cannot represent.")
\end{lstlisting}
The output is worth reading carefully, because the interesting column is the one
that does \emph{not} change:

\smallskip
\noindent\begin{tabular}{@{}L{3.4cm}rrrr@{}}
\toprule
\textbf{Agent} & \textbf{Total} & \textbf{Breaches} & \textbf{On-call} &
\textbf{Escalations} \\
\midrule
Simple reflex & 259 & 136 & 123 & 30 \\
Model-based   & 156 & 136 &  20 &  5 \\
\bottomrule
\end{tabular}

\smallskip
The breach cost is \emph{identical}: $136$ either way. The simple reflex agent's
extra twenty-five escalations bought nothing at all, because past the third one
Ayesha was too tired to fix anything --- and the sleep was lost regardless. The
total falls by 40\% purely by not doing the useless work, and the escalations fall
sixfold.

Now change \texttt{WAKE\_LIMIT} to $40$ and re-run. The two agents converge,
because with an unlimited budget the number being tracked no longer constrains
anything. That is the test of whether internal state is earning its keep: state is
worth maintaining exactly when it represents a limit the environment imposes and
the percepts do not report.
\end{lab}

%---------------------------------------------------------------------
\begin{checkpoint}
\begin{tightnum}
  \item Write a PEAS description for an agent that decides which of a day's pull
        requests to review first. Make sure the performance measure describes the
        state of the repository, not the activity of the agent.
  \item A vending machine dispenses a drink when the right coins are inserted.
        Classify its environment on all six properties, and say which of the five
        agent architectures it needs.
  \item An agent escalates a ticket on strong evidence and the escalation turns
        out to have been unnecessary. Was it rational? Justify your answer in one
        sentence using the definition in this chapter.
\end{tightnum}
\end{checkpoint}

%---------------------------------------------------------------------
\begin{chaptersummary}
\begin{tight}
  \item An agent perceives through sensors and acts through actuators. The
        \term{agent function} maps every percept sequence to an action; the
        \term{agent program} is a bounded implementation of it. AI is the search
        for good programs, because the functions are too large to tabulate.
  \item Specify every task as \term{PEAS} before building anything. The
        performance measure must describe the state of the environment, and must
        be written so that two people would score the same run identically.
  \item Classify the environment on six properties: observable, deterministic,
        episodic, static, discrete, single-agent. Each one forces a design
        decision, and partial observability is the most consequential --- it is
        what makes internal state necessary.
  \item There are five agent architectures: simple reflex, model-based reflex,
        goal-based, utility-based and learning. Use the simplest that can do the
        job. This course builds the middle three; \chref{learning} signposts the
        fifth.
  \item Rationality is maximising \emph{expected} performance given the percepts
        so far. It is not omniscience, not perfection, and it positively
        recommends gathering information when information is cheap.
\end{tight}
\end{chaptersummary}

\begin{reviewq}
  \item \textbf{(MCQ)} The agent function maps: (a)~states to actions (b)~percept
        sequences to actions (c)~actions to utilities (d)~percepts to goals.
  \item \textbf{(MCQ)} An agent must keep internal state principally because the
        environment is: (a)~stochastic (b)~dynamic (c)~partially observable
        (d)~multi-agent.
  \item \textbf{(MCQ)} Which is \emph{not} part of a PEAS description?
        (a)~performance measure (b)~environment (c)~algorithm (d)~sensors.
  \item \textbf{(MCQ)} Chess played against a clock is: (a)~fully observable,
        deterministic, static (b)~fully observable, deterministic, dynamic
        (c)~partially observable, stochastic, dynamic (d)~fully observable,
        stochastic, static.
  \item \textbf{(Short)} Explain the difference between rationality and
        omniscience, and say why the distinction matters when reviewing an
        incident.
  \item \textbf{(Short)} Give one task for which a simple reflex agent is the
        right choice, and one superficially similar task for which it is not.
  \item \textbf{(Short)} Why must the performance measure describe the environment
        rather than the agent? Give a measure that fails this test and say how it
        would be exploited.
  \item \textbf{(Applied)} Take the sprint scheduler of Table~\ref{tab:peas}. Write its
        full PEAS description, classify the environment on all six properties, and
        state which agent architecture the classification forces. Then name the
        single simplification that \chref{csp} makes in order to solve it at all,
        and what is lost by making it.
\end{reviewq}
